\documentclass[11pt,a4paper]{article}
\usepackage[T1]{fontenc}\usepackage[utf8]{inputenc}\usepackage[italian]{babel}\usepackage{lmodern,microtype}
\usepackage[a4paper,margin=2.4cm,headheight=15pt]{geometry}\usepackage{enumitem,booktabs,longtable,array}\usepackage{xcolor,hyperref,xurl,fancyhdr,titlesec,listings,framed}
\definecolor{ddtadark}{RGB}{28,38,52}\definecolor{ddtagray}{RGB}{242,244,247}\definecolor{ddtablue}{RGB}{42,88,135}\definecolor{ddtagreen}{RGB}{48,111,78}\definecolor{ddtaorange}{RGB}{148,91,25}
\hypersetup{colorlinks=true,linkcolor=ddtablue,urlcolor=ddtablue,pdftitle={DDTA Thesis Chapters 2-4 Closure Record},pdfauthor={Documentation-Driven Threat Analysis research workspace}}
\pagestyle{fancy}\fancyhf{}\lhead{DDTA thesis closure}\rhead{Chapters 2--4}\cfoot{\thepage}
\titleformat{\section}{\Large\bfseries\color{ddtadark}}{\thesection}{0.7em}{}\titleformat{\subsection}{\large\bfseries\color{ddtadark}}{\thesubsection}{0.7em}{}
\setlist[itemize]{leftmargin=1.55em,itemsep=0.25em,topsep=0.35em}\setlength{\parskip}{0.35em}\setlength{\parindent}{0pt}\setlength{\emergencystretch}{2em}
\lstset{basicstyle=\ttfamily\small,columns=fullflexible,keepspaces=true,frame=single,framerule=0.3pt,rulecolor=\color{ddtagray},backgroundcolor=\color{ddtagray},breaklines=true}
\newenvironment{closedbox}[1]{\def\FrameCommand{\fcolorbox{ddtagreen}{ddtagray}}\MakeFramed{\advance\hsize-2\width\FrameRestore}\noindent\textbf{#1}\par\smallskip}{\endMakeFramed}
\newenvironment{reopenbox}[1]{\def\FrameCommand{\fcolorbox{ddtaorange}{ddtagray}}\MakeFramed{\advance\hsize-2\width\FrameRestore}\noindent\textbf{#1}\par\smallskip}{\endMakeFramed}
\begin{document}
\begin{titlepage}\centering\vspace*{2cm}{\Large Documentation-Driven Threat Analysis\par}\vspace{0.8cm}{\Huge\bfseries Thesis Chapters 2--4\par}\vspace{0.3cm}{\LARGE Closure and reopen record\par}\vspace{1cm}{\large Freeze del background, research gap e documentation metamodel prima della Base Analysis\par}\vfill\begin{minipage}{0.84\textwidth}\small
\textbf{Baseline precedente alla finalizzazione.} \nolinkurl{ec6c0107e0d5c0460d2afcbe12dae6f40bc5e6c5}.\\[0.5em]
\textbf{Significato di FINAL.} Stabile per lo scope corrente della tesi, non immutabile. Una riapertura richiede un criterio esplicito e un controesempio materiale.\end{minipage}\vfill{\large 15 agosto 2026\par}\end{titlepage}
\tableofcontents\newpage
\section{Scopo del record}
Questo record stabilisce un freeze di ricerca per il blocco della tesi che motiva e definisce il layer documentale governato prima di iniziare la Base Analysis. La chiusura impedisce che i capitoli già stabilizzati vengano continuamente riscritti per accomodare dettagli appartenenti ai layer analitici o di tooling.

\section{Capitolo 2 -- Background}
\begin{closedbox}{CLOSED / FINAL}
Il Capitolo 2 è semanticamente chiuso. Non richiede una nuova elaborazione concettuale prima della Base Analysis.
\end{closedbox}
Conserva le distinzioni secondo cui threat analysis può iniziare prima dell'implementazione, documentazione e modello analizzabile non sono sinonimi, metodologie diverse possono richiedere viewpoint diversi, Finding/minaccia/mitigazione/SecurityRequirement sono artefatti distinti, output automatici o LLM restano candidati fino a review, e traceability/provenance/change-awareness influenzano la validità analitica.
\begin{reopenbox}{Reopen clause -- Capitolo 2}
Riaprire solo per errore fattuale/bibliografico materiale, claim materialmente non supportato o falsificazione di una distinzione fondativa usata dalla tesi. Un nuovo BAE o una nuova relazione non costituiscono da soli un trigger.
\end{reopenbox}

\section{Capitolo 3 -- Stato dell'arte e research gap}
\begin{closedbox}{CLOSED / FINAL dopo correzione editoriale di coerenza}
G1--G4 e lo scope della tesi restano invariati. La finalizzazione corregge soltanto un forward reference ormai stale.
\end{closedbox}
Il research gap resta composto da: documentazione portable-by-construction verso rappresentazione analizzabile revisionata; common core methodology-neutral; risultati method-specific verso finding revisionati e requisiti di sicurezza governati; change-aware traceability e re-analysis.
\begin{reopenbox}{Reopen clause -- Capitolo 3}
Riaprire solo se il research gap o lo scope valutato vengono materialmente falsificati/modificati, oppure se prior work viene trovato a risolvere una discontinuità dichiarata nello stesso senso end-to-end valutato dalla tesi.
\end{reopenbox}

\section{Capitolo 4 -- Documentation metamodel and authoring rules}
\begin{closedbox}{CLOSED / FINAL dopo consolidamento S1.5 + S2}
Il Capitolo 4 incorpora definitivamente il documentation metamodel fino a SecurityRequirement.
\end{closedbox}
Gerarchia di tipo:
\begin{lstlisting}
GovernedDocument
    |
    +-- MacroRequirement
    +-- Decision
    `-- Requirement [abstract]
            normativeClause : NormativeClause [1..*]
            |
            +-- FunctionalRequirement
            `-- SpecializedRequirement [abstract]
                    `-- SecurityRequirement
\end{lstlisting}
Percorso concreto di authoring/ownership:
\begin{lstlisting}
Project problem framing [method precondition]
    -> MacroRequirement
    -> Decision
    -> FunctionalRequirement
    -> SecurityRequirement
       (concrete SpecializedRequirement)
\end{lstlisting}
\subsection{Risultati consolidati}
\begin{itemize}
\item MR e single macro responsibility;
\item Decision come significant commitment sotto un solo MR;
\item FR come Requirement operativo indipendentemente assessable sotto una sola Decision;
\item \texttt{Requirement [abstract]} come obbligazione normativa governata comune;
\item \texttt{Requirement.normativeClause : NormativeClause [1..*]};
\item coherent-unit e split-on-independence;
\item \texttt{NormativeObligation} come metaclasse L1 separata REJECTED;
\item SpecializedRequirement come strengthening normativo di esattamente un FR;
\item composizione congiuntiva, removal test, realization separation e boundary con functional correctness;
\item \texttt{SecurityRequirement IS-A SpecializedRequirement};
\item \texttt{protectedSecurityProperty : SecurityProperty [1]};
\item failure-mode explicitness nelle normative clauses;
\item cause neutrality e single governing security property;
\item Attack/Finding/Risk/Control/evidence esclusi dalla semantica core del SecurityRequirement.
\end{itemize}
\subsection{Elementi ancora aperti ma non appartenenti alla closure del Capitolo 4}
Lifecycle/history/change-event completo; esatta rappresentazione L2 di NormativeClause; tassonomia SecurityProperty; possibile identità strutturale futura di SecurityFailureMode; Base Analysis/BAE; AnalysisRecord/Finding; STRIDE/STRIDE-AI; Risk/Control/evidence; provenance analitica strutturale.
\begin{reopenbox}{Reopen clause -- Capitolo 4}
Riaprire solo se una fase successiva produce un failure semantico concreto e ricorrente non rappresentabile senza cambiare il contratto documentale governato. Il controesempio deve identificare il minimo owning layer e attivare regressione dei corpus accettati.
\end{reopenbox}
Non riaprono da soli il capitolo: un nuovo tipo/relazione BAE, una proiezione grafica, un campo privato di plugin, un sinonimo suggerito dal tool, un candidato LLM, una nuova tecnica di implementazione/control.

\section{Freeze rule}
Dopo il commit di finalizzazione:
\begin{lstlisting}
Chapter 2  CLOSED / FINAL
Chapter 3  CLOSED / FINAL
Chapter 4  CLOSED / FINAL
\end{lstlisting}
Il lavoro successivo deve referenziare questi capitoli invece di retrofittarli continuamente. Correzioni editoriali restano possibili; la riapertura semantica richiede i criteri sopra.

\section{Prossimo confine di ricerca}
Il prossimo blocco autorizzato è Base Analysis. Il primo microstep è definire responsabilità e confine della rappresentazione analizzabile methodology-neutral prima di fissare una tassonomia BAE o applicare un threat method.
\end{document}
